\section{Radial-velocity measurement and binary classification}
\label{sec:methods}

\subsection{Cross-correlation function}
\label{sec:ccf}

Radial velocities were measured using the one-dimensional
cross-correlation function (CCF) technique of \citet{Zucker1994}, in the
same form as adopted for the parallel Galactic WC survey of
\citet{Dsilva2020}. For two flux arrays $f_1$ and $f_2$ defined on a
common logarithmic wavelength grid of length $N$, the CCF is the
normalised dot product
\begin{equation}
\rho(s) \;=\; \frac{1}{N\,\sigma_{f_1}\,\sigma_{f_2}}
              \sum_{i=1}^{N} f_1(\lambda_i)\,f_2(\lambda_i + s),
\end{equation}
where $s$ is the velocity shift, and the velocity shift of $f_2$ relative
to $f_1$ is recovered from the position of the CCF maximum. For each star,
a high signal-to-noise template spectrum was constructed by first
computing the CCF of all epochs against an arbitrary reference epoch to
determine relative velocity shifts, correcting each epoch to a common rest
frame, and then co-adding the corrected spectra with signal-to-noise
squared weighting. This iterative procedure removes the dependence on the
choice of reference epoch.

The CCF was computed independently for each of the \NlinesAll emission
lines listed in Table~\ref{tab:emission_lines}, spanning the full
X-SHOOTER wavelength range from the ultraviolet
(\ion{O}{v}\,$\lambda\lambda 3100$--$3175$) to the near-infrared
(\ion{C}{iv}\,$\lambda 20842$). Each line is associated with a default
integration window in wavelength (Table~\ref{tab:emission_lines}) and a
default CCF fit fraction $f_\mathrm{fit} = \fitFracDefault$, which
controls the fraction of the CCF peak used for parabolic fitting: the
fit region extends to the wavelengths at which
$\rho(s) = f_\mathrm{fit}\,\rho_{\max}$, and a parabola of the form
$\rho(s) \simeq a s^2 + b s + c$ is fitted to the points above this
threshold. The radial velocity is then recovered from the parabola apex,
$v = -b/(2a)$, and the formal per-epoch uncertainty
$\sigma_\mathrm{RV}$ is derived from the curvature of the CCF peak
following \citet{Zucker2003},
\begin{equation}
\sigma_\mathrm{RV}
   \;=\; \left[\,
     -\,N \,\frac{\rho''(\hat s)\,\rho(\hat s)}{1 - \rho(\hat s)^2}
   \,\right]^{-1/2},
\end{equation}
where $\hat s$ is the position of the CCF maximum and $\rho''$ is the
second derivative of the CCF at the peak.

For a small fraction of star--line combinations the default settings
yielded a poorly converged parabolic fit -- typically because of edge
effects from the chosen wavelength range, or a noise spike close to the
peak -- and the wavelength range and/or fit fraction were adjusted on a
per-star, per-line basis. All such overrides are documented in the public
CCF settings file released alongside this paper, and their net effect on
the RV measurements is negligible compared to the per-epoch
$\sigma_\mathrm{RV}$.

% Emission line table
\begin{table}
\caption{Emission lines used for cross-correlation analysis.}
\label{tab:emission_lines}
\centering
\begin{tabular}{lcc}
\hline\hline
Line ID & $\lambda_\mathrm{min}$ [nm] & $\lambda_\mathrm{max}$ [nm] \\
\hline
\ion{O}{v}\,$\lambda\lambda 3100$--$3175$   & 310.0 & 317.5 \\
\ion{O}{iv}\,$\lambda\lambda 3350$--$3480$  & 335.0 & 348.0 \\
\ion{C}{iv}\,$\lambda\lambda 3650$--$3900$  & 365.0 & 390.0 \\
\ion{He}{ii}\,$\lambda 4686$                & 456.0 & 480.0 \\
\ion{O}{vi}\,$\lambda\lambda 5210$--$5340$  & 521.0 & 533.5 \\
\ion{He}{ii}\,$\lambda 5412$ \& \ion{C}{iv}\,$\lambda 5471$ & 535.0 & 554.0 \\
\ion{C}{iv}\,$\lambda\lambda 5808$--$5812$  & 570.0 & 588.0 \\
\ion{C}{iii}\,$\lambda\lambda 6700$--$6800$ & 666.5 & 684.0 \\
\ion{C}{iv}\,$\lambda 7063$                 & 697.0 & 714.0 \\
\ion{C}{iv}\,$\lambda 17396$                & 1710.0 & 1763.0 \\
\ion{C}{iv}\,$\lambda 20842$                & 2050.0 & 2100.0 \\
\hline
\end{tabular}
\end{table}

\subsection{Choice of primary RV line}
\label{sec:primary_line}

At the outset of the analysis we made no \emph{a priori} commitment to
which of the \NlinesAll emission lines would yield the most reliable
radial velocity. We therefore measured the CCF for all of them, and
selected the primary RV line based on the physical and empirical
considerations set out below.

\paragraph{NIR lines and telluric H$_2$O contamination.}
The two near-infrared carbon features in our line list,
\ion{C}{iv}\,$\lambda 17396$ and \ion{C}{iv}\,$\lambda 20842$, fall on
the red wing of the strong $2\nu_2 + \nu_3$ / $\nu_1 + \nu_3$ H$_2$O
combination band system between 1.75 and 2.00\,$\mu$m and on the blue
edge of the 2.0\,$\mu$m H$_2$O band, respectively. Atmospheric
transmission at Cerro Paranal in these regions drops below 70--80\%
even at moderate precipitable water vapour
($\mathrm{PWV} \simeq 2.5$\,mm), and telluric correction with
\texttt{molecfit} \citep{Smette2015,Kausch2015} leaves residuals at the
few-percent level. On the broad ($\Delta v \sim 10^3$\,\kms) WC
emission profiles such residuals deform the line shape and bias the
CCF centroid by velocities comparable to or larger than the binary
detection threshold $\drvthresh$\,\kms. Spectroscopic monitoring of
WC9 stars by \citet{Williams2000} corroborates the unreliability of
NIR CCF centroids in this regime, and we therefore exclude the two NIR
lines from the primary classification.

\paragraph{\ion{He}{ii} variability and blending.}
The \ion{He}{ii}\,$\lambda 4686$ line forms higher in the accelerating
WC wind than the deeper \primaryline doublet
\citep{Hillier1989,Crowther2007}, making its profile more sensitive to
wind clumping, co-rotating interaction regions, and excess emission
from wind--wind collision in binaries. Stochastic profile variability
at the few-percent level on hour-to-day timescales is well documented
\citep{LepineMoffat1999}, and \citet{Shenar2019} explicitly recommend
avoiding \ion{He}{ii}\,$\lambda 4686$ for radial-velocity measurements
when alternatives exist. The \ion{He}{ii}\,$\lambda 5412$ line is in
addition blended with the \ion{C}{iv}\,$\lambda 5471$ feature in WC
spectra (Table~\ref{tab:emission_lines}), so that wind variability and
line confusion both contribute to a less stable centroid.

\paragraph{The \primaryline doublet -- primary RV line.}
The \primaryline doublet is the strongest and deepest-formed optical
emission feature in the WC spectrum and yields the most stable line
centroid in WC atmosphere modelling. We therefore adopt \primaryline
as the primary RV diagnostic for the binary classification of
Sect.~\ref{sec:criteria}. This choice also provides direct
methodological continuity with the precursor LMC WC binary survey of
\citet{Bartzakos2001}, which used the same line, and with the Galactic
WC survey of \citet{Dsilva2020,Dsilva2023}.

The remaining lines in Table~\ref{tab:emission_lines} -- the
\ion{O}{v}, \ion{O}{iv}, \ion{O}{vi}, \ion{C}{iv}\,$\lambda 3650$,
\ion{C}{iii}\,$\lambda 6700$ and \ion{C}{iv}\,$\lambda 7063$ features
-- are retained as cross-checks of the per-star RV time series
(Sect.~\ref{sec:line_agreement}), but are not used to flag binaries.

\subsection{Signal-to-noise requirements}
\label{sec:snr}

To ensure reliable radial-velocity measurements, we established minimum
signal-to-noise requirements per emission line through simulations. We
generated 30 synthetic single-lined binary (SB1) spectra using a WC4
template with varying S/N levels, measured the radial-velocity scatter for
each line, and required $\sigma_\mathrm{RV} < 3$\,\kms as the acceptance
threshold.

\subsection{Line-agreement cross-check}
\label{sec:line_agreement}

The choice of \primaryline as the primary RV line is supported by an
independent, data-driven cross-check across all \NlinesAll lines.

For each pair of emission lines $(\ell, \ell')$ we computed the Pearson
correlation coefficient $r_{\ell\ell'}$ of the per-star peak-to-peak
$\Delta\mathrm{RV}_{\max}$ values, restricted to the stars in which both
lines yielded valid measurements. To weight the correlation by sample
size, we define the agreement score per line $\ell$ as
\begin{equation}
S_\ell \;=\; \frac{\sum_{\ell' \neq \ell}
                    w_{\ell\ell'}\,r_{\ell\ell'}}
                  {\sum_{\ell' \neq \ell} 1},
\qquad
w_{\ell\ell'} \;=\; \frac{n_{\ell\ell'}}{\Nstars},
\label{eq:agreement_score}
\end{equation}
where $n_{\ell\ell'}$ is the number of stars with valid
$\Delta\mathrm{RV}_{\max}$ from both lines; we require
$n_{\ell\ell'} \geq \Nstarsmincorr$ for a pair to enter the average.
The line that is most ``agreeable'' with the rest of the sample has
the highest $S_\ell$.

The \primaryline doublet emerges as the highest-scoring line. As an
independent check, we computed the per-line equivalent width on the
same dataset and found a positive correlation between $S_\ell$ and the
mean per-epoch line strength: the strongest, most isolated WC emission
feature is therefore also the most kinematically self-consistent
across the sample. This is an empirical confirmation of the
depth-of-formation argument set out in Sect.~\ref{sec:primary_line}.
The agreement-score ranking and the score-vs-strength relation are
shown in Sect.~\ref{sec:results} (Fig.~\ref{fig:agreement}).

\subsection{$\Delta\mathrm{RV}$ threshold derivation}
\label{sec:threshold_derivation}

The binary detection threshold $\drvthresh$\,\kms used in
Sect.~\ref{sec:criteria} was derived empirically from the cumulative
observed binary fraction as a function of threshold.

\paragraph{Parametric Gaussian model.}
We model the cumulative binary fraction above threshold $T$ as the
survival function of a two-component Gaussian mixture, in which the
peak-to-peak $\Delta\mathrm{RV}_{\max}$ distributions of the
single-star and unresolved-binary populations are both approximated as
zero-centred normals,
\begin{equation}
f_\mathrm{bin}(>T)
  \;=\; (1 - f_\mathrm{bin})\,\Phi(T/\sigma_s)
        \;+\; f_\mathrm{bin}\,\Phi(T/\sigma_b),
\label{eq:gauss_threshold}
\end{equation}
where $\Phi$ is the survival function (i.e.\ $1 - \mathrm{CDF}$) of a
standard normal, $\sigma_s$ is the single-star peak-to-peak RV
scatter, and $\sigma_b$ is the corresponding scatter of the binary
population. The free parameters $(\sigma_s, \sigma_b, f_\mathrm{bin})$
are fitted to the observed $f_\mathrm{bin}(>T)$ vs $T$ curve by
non-linear least squares. The fitted parameters are
$\sigma_s = \sigmaSingleFit$\,\kms,
$\sigma_b = \sigmaBinaryFit$\,\kms, and
$f_\mathrm{bin} = \fbinAnalytic$, and the derived binary-detection
threshold is $T = \drvthresh$\,\kms, the value adopted throughout this
work.

\paragraph{Threshold value.}
The adopted threshold is the inflection of the modelled
$f_\mathrm{bin}(>T)$ curve: below it, the cumulative detection rate is
dominated by single-star wind variability, while above it, it is
dominated by genuine binaries. We adopt $\drvthresh$\,\kms as the
canonical binary-detection threshold (Sect.~\ref{sec:criteria}),
subject to the caveat below.

\paragraph{Caveat: the Gaussian assumption is approximate.}
Equation~\eqref{eq:gauss_threshold} assumes that both the single-star
and the unresolved-binary $\Delta\mathrm{RV}_{\max}$ distributions are
Gaussian. WR wind variability is, however, intrinsically stochastic
and asymmetric: clumping events, co-rotating interaction regions, and
line-profile excursions all produce non-Gaussian wings that the
analytic model of Eq.~\eqref{eq:gauss_threshold} cannot capture. We
have constructed a diagnostic tool that compares the observed
single-star $\Delta\mathrm{RV}_{\max}$ distribution to the best-fit
Gaussian via residual and second-derivative tests, and the diagnostic
confirms that the Gaussian is a useful first approximation but is not
exact in the wings. The bias-correction analysis described in
Sect.~\ref{sec:bias} is consequently formulated as a non-parametric
multinomial likelihood comparison (Eq.~\ref{eq:multinomial}) that does
not assume Gaussianity, and the threshold-curve cross-check in
Sect.~\ref{sec:threshold_fit} uses a Physics-based simulation rather
than an analytic Gaussian. A detailed treatment of the non-Gaussian
wind-variability distribution is deferred to a companion paper.

\subsection{Binary classification criteria}
\label{sec:criteria}

A star is classified as a spectroscopic binary if \emph{both} of the
following criteria are satisfied \citep[following][]{Bartzakos2001,
Dsilva2023}:

\begin{enumerate}
\item The peak-to-peak radial-velocity difference between any pair of
  epochs exceeds $\Delta\mathrm{RV}_\mathrm{max} > \drvthresh$\,\kms.
\item The detection is statistically significant:
  $\Delta\mathrm{RV} - \sigthresh\sigma > 0$, where
  $\sigma = \sqrt{\sigma_1^2 + \sigma_2^2}$ is the combined formal
  uncertainty of the epoch pair.
\end{enumerate}

The classification is performed in two stages. First, the epoch pair with
the maximum radial-velocity separation is evaluated. If this pair does not
satisfy both criteria, all remaining epoch pairs are scanned. The binary
decision is based exclusively on the \primaryline emission line, the
strongest and most isolated WC spectral feature
(Sect.~\ref{sec:primary_line}).

\subsection{Threshold-curve fit}
\label{sec:threshold_fit}

As a complementary check to the grid-based likelihood analysis described in
Sect.~\ref{sec:bias}, the cumulative observed binary fraction is fitted as
a function of the detection threshold $T$ on $\Delta\mathrm{RV}_\mathrm{max}$.
We adopt the Physics-based simulation: orbital parameters (period,
eccentricity, mass ratio, inclination) are drawn from the same priors
as in Sect.~\ref{sec:bias}, and the radial-velocity curve of each
synthetic binary is computed via Kepler's third law (Eq.~\ref{eq:K1})
and sampled at the observed cadence of each target. The free parameters
are the intrinsic binary fraction $f_\mathrm{bin}$, the single-star
scatter $\sigma_\mathrm{intrinsic}$, and the per-epoch measurement
uncertainty $\sigma_\mathrm{meas}$. Single stars are assigned intrinsic scatter at scale
$\sigma_\mathrm{intrinsic}$ and per-epoch noise at scale
$\sigma_\mathrm{meas}$, and the same two-criterion detection of
Sect.~\ref{sec:criteria} is applied. The cumulative binary fraction
$N_\mathrm{bin}(>T)/\Ntotal$ is computed across a grid of thresholds $T$
and compared to the observed cumulative curve to extract the best-fit
parameters.

\subsection{Spectroscopic morphology investigation}
\label{sec:spec_morphology}

In addition to the RV-based classification of Sect.~\ref{sec:methods}, each
star is visually inspected through a diagnostic-line viewer that groups the
spectrum into curated zoom panels covering interstellar features, telluric
bands, hydrogen Balmer and \ion{He}{i}/\ion{He}{ii} OB-type photospheric
absorption, the WR \ion{He}{ii}\,$\lambda 4686$ emission, the
\primaryline emission complex, and the
\ion{O}{v}/\ion{O}{vi} lines that distinguish WC from WO subtypes. Stars
showing morphological signatures inconsistent with a single WC/WO spectrum
-- for example, weak photospheric absorption components attributable to an
OB companion, or anomalous emission-line profiles -- are flagged as
candidate binaries. If any additional binaries are flagged through this
morphological inspection, the binary count is updated and the Monte-Carlo
bias correction is re-run with the revised \Nbin and \Ntotal. This
inspection adds \NbinSpec spectroscopically flagged binaries to the
RV-based count, raising the total to \NbinFinal/\Ntotal and the observed
binary fraction to $\fbinFinal$.
